% Methods and Results sections, IEEEtran (\input into main.tex).
% Requires: amsmath, booktabs, graphicx, siunitx not needed (units written literally).
% Figure assets referenced here (fig_block_diagram, fig_setup, fig_offsets) are not yet
% drawn; each figure environment carries a visible placeholder so the file compiles.

\section{Methods}

\subsection{Landing as an Image-Plane Regulation Problem}

A quadrotor descending onto a deck that is itself moving cannot rely on a position
estimate expressed in a world frame, because any error in that estimate accumulates over
the descent and is unobservable to the vehicle. The relative quantity that matters for
landing, namely the horizontal offset between the vehicle and the deck, is measured
directly by a downward-facing camera that observes a fiducial marker fixed to the deck.
We therefore regulate that offset in the image plane and never form a world-frame estimate
of the deck.

The marker detection is smoothed with an exponential moving average filter, which
suppresses the pixel-level jitter of individual detections without adding a state estimator
between the camera and the control law. Writing $\bar{\mathbf{p}} = (\bar{u}, \bar{v})^{\top}$
for the smoothed marker center, the image-plane error is given by
%
\begin{equation}
  \mathbf{e} = \bar{\mathbf{p}} - \mathbf{c},
  \label{eq:error}
\end{equation}
%
where $\mathbf{c} = (u_{c}, v_{c})^{\top}$ is the principal point, taken at the image
center. The error is normalized per axis by the corresponding half-image dimension,
$\hat{e}_{u} = e_{u}/u_{c}$ and $\hat{e}_{v} = e_{v}/v_{c}$, so that the control gains are
independent of image resolution.

The horizontal velocity command is proportional to the normalized error and damped by the
measured horizontal velocity of the vehicle,
%
\begin{align}
  v_{x} &= \mathrm{sat}\!\left(-k_{p}\,\hat{e}_{v} - k_{d}\,\tilde{v}_{x},\ \bar{v}_{xy}\right), \label{eq:vx}\\
  v_{y} &= \mathrm{sat}\!\left(-k_{p}\,\hat{e}_{u} - k_{d}\,\tilde{v}_{y},\ \bar{v}_{xy}\right), \label{eq:vy}
\end{align}
%
where $k_{p}$ is the proportional gain, $k_{d}$ is the damping gain,
$\tilde{v}_{x}$ and $\tilde{v}_{y}$ are the horizontal velocities reported by the flight
controller, $\bar{v}_{xy}$ is the horizontal speed limit, and
$\mathrm{sat}(x, a) = \min(\max(x, -a), a)$. The image rows are aligned with the forward
axis of the vehicle and the image columns with its lateral axis, so the row component of
the error drives the forward velocity and the column component drives the lateral velocity.
The damping term uses the measured vehicle velocity rather than a numerical derivative of
the error, which keeps the command free of differentiated detection noise and, on a moving
deck, resists the standing velocity error that a purely proportional law would leave.

Descent is gated on alignment. The vertical velocity command is
%
\begin{equation}
  v_{z} =
  \begin{cases}
    -v_{f}, & h < h_{f},\\
    -\bar{v}_{z}, & h \geq h_{f} \ \text{and}\ \lVert \mathbf{e} \rVert < \rho_{g},\\
    0, & \text{otherwise},
  \end{cases}
  \label{eq:vz}
\end{equation}
%
where $h$ is the estimated altitude above the deck, $h_{f}$ is the flare altitude,
$\bar{v}_{z}$ is the nominal descent speed, $v_{f}$ is the reduced descent speed used
during flare, and $\rho_{g}$ is the alignment radius. The gate is evaluated on the raw
pixel error of \eqref{eq:error} rather than on the normalized error, so $\rho_{g}$ is
expressed in pixels. Holding altitude until the vehicle is aligned prevents the descent
from consuming the control authority that lateral correction requires, and the reduced
speed below $h_{f}$ limits the impact energy at contact. Parameter values are listed in
Table~\ref{tab:params}.

Velocity setpoints are streamed to the flight controller at $f_{c} = 50$~Hz, whereas
marker detections arrive at $f_{d} = 8$~Hz. Most control cycles therefore execute without
a new measurement, and the smoothed detection is held between frames. When no detection
has been received for longer than $\tau_{h}$, the last command produced from a fresh
detection is reissued at half amplitude until detections resume, which keeps the vehicle
moving with the deck instead of arresting over a marker it can no longer see.

\subsection{Altitude Estimation}

Neither available altitude source is sufficient alone. The barometer is continuous but
drifts, and the rangefinder is accurate near the ground but noisy. The two are blended as
%
\begin{equation}
  h = (1 - \kappa)\, h_{b} + \kappa\, h_{r},
  \qquad
  \kappa = \frac{\Delta t}{\tau + \Delta t},
  \label{eq:alt}
\end{equation}
%
where $h_{b}$ is the barometric altitude, $h_{r}$ is the rangefinder altitude, $\Delta t$
is the control period, and $\tau$ is the blending time constant. When no rangefinder
sample is available, $h = h_{b}$. With the values of Table~\ref{tab:params} the estimate
follows the barometer within one control cycle and the rangefinder enters as a small
persistent correction. The altitude estimate is read only by the flare condition of
\eqref{eq:vz}.

%% TODO(author): \eqref{eq:alt} is the blend the code actually implements, which is a
%% static weighting rather than the complementary filter the source comment describes.
%% Decide whether to correct the implementation and re-run before submission, or to keep
%% the present description of the as-flown system.

\subsection{Baseline Controller}

The comparison isolates the contribution of the visual measurement. The
\emph{vision-servoed controller} is the system of \eqref{eq:error}--\eqref{eq:vz}. The
\emph{open-loop controller} commands a constant descent at $\bar{v}_{z}$ with no lateral
velocity and no flare, and is the dead-reckoning strategy available when the deck is not
observed. The two controllers share the vehicle, the flight-controller interface, the
command rate, and the initial hover pose, and differ only in whether the marker
measurement enters the command.

\subsection{Experimental Setup}

Landings were performed on a 450~mm quadrotor running an open-source flight
stack\footnote{PX4 v1.14.} and carrying a downward-facing global-shutter camera at
$640 \times 480$ pixels. The camera detected an AprilTag 36h11 marker of 160~mm side
length mounted at the center of the deck of a wheeled cart, which travelled at a constant
0.3~m/s. The deck was circular with a radius of 300~mm, which defines the physical
tolerance the landing must meet. Ground truth was recorded by a motion-capture system
(OptiTrack) at 120~Hz, which tracked both the vehicle and the deck, and all quantities
reported in Section~\ref{sec:results} are derived from that record rather than from
onboard estimates. Fig.~\ref{fig:setup} shows the apparatus.

Each controller was evaluated over $N = 12$ landings, and the two sets were analyzed as 12
paired observations. Every landing began from a hover at the same altitude and lateral
offset from the moving deck, and ended at contact.

%% TODO(author): state the pairing scheme (alternating? blocked? randomized?) and the
%% initial hover altitude/offset --- neither is recorded in the lab log.

\subsection{Metrics and Statistical Analysis}

The primary metric is the \emph{touchdown offset}, defined as the horizontal distance
between the point of first contact and the deck center at the same instant, both taken
from the motion-capture record. A landing is counted as successful when the touchdown
offset is smaller than the 300~mm deck radius. The \emph{descent duration} is the elapsed
time from the start of the descent to first contact.

Touchdown offset is summarized as mean $\pm$ standard deviation across trials, which
reports the achieved accuracy and its repeatability, and the difference between
controllers is summarized as the paired mean difference with a 95\% confidence interval,
which reports the size and the precision of the improvement. The paired difference was
tested with a two-sided paired $t$-test at a significance level of 0.05. Inference was
performed on the primary metric only. Descent duration and landing success are reported
descriptively, so no correction for multiplicity is required.

\begin{table}[!t]
\caption{Controller parameters. Values apply to the vision-servoed controller. The
open-loop controller uses only the nominal descent speed $\bar{v}_{z}$.}
\label{tab:params}
\centering
\footnotesize
\setlength{\tabcolsep}{3pt}
\begin{tabular}{llr}
\toprule
Symbol & Description & Value \\
\midrule
$f_{c}$        & Velocity setpoint rate (Hz)                          & 50 \\
$f_{d}$        & Marker detection rate (Hz)                           & 8 \\
$k_{p}$        & Proportional gain on normalized image error (m/s)    & 0.9 \\
$k_{d}$        & Damping gain on measured horizontal velocity         & 0.15 \\
$\bar{v}_{xy}$ & Horizontal speed limit (m/s)                         & 1.2 \\
$\bar{v}_{z}$  & Nominal descent speed (m/s)                          & 0.5 \\
$\rho_{g}$     & Alignment radius for the descent gate (px)           & 40 \\
$h_{f}$        & Flare altitude above the deck (m)                    & 0.35 \\
$v_{f}$        & Descent speed during flare (m/s)                     & 0.25 \\
$\tau$         & Altitude blending time constant (s)                  & 0.8 \\
$\tau_{h}$     & Detection timeout before the half-amplitude hold (s) & 0.4 \\
\bottomrule
\end{tabular}
\end{table}

\begin{figure}[!t]
\centering
% \includegraphics[width=\columnwidth]{fig_block_diagram.pdf}
\fbox{\parbox[c][22mm][c]{0.92\columnwidth}{\centering\footnotesize
Placeholder: linear left-to-right block diagram --- camera $\rightarrow$ marker detection
$\rightarrow$ smoothing $\rightarrow$ image-plane error $\rightarrow$ velocity law
$\rightarrow$ flight controller, with the altitude blend feeding the flare condition.}}
\caption{System architecture of the vision-servoed landing controller. Marker detections
from the downward camera are smoothed and differenced against the principal point to form
the image-plane error, which the velocity law converts into horizontal and vertical
setpoints for the flight controller. The blended altitude estimate enters only the flare
condition. Feedback through the physical world is implicit.}
\label{fig:block}
\end{figure}

\begin{figure}[!t]
\centering
% \includegraphics[width=\columnwidth]{fig_setup.pdf}
\fbox{\parbox[c][22mm][c]{0.92\columnwidth}{\centering\footnotesize
Placeholder: annotated photograph of the apparatus with callouts on the quadrotor, the
downward camera, the AprilTag marker, the cart deck, and the motion-capture cameras.}}
\caption{Experimental apparatus. Callouts identify the 450~mm quadrotor, its
downward-facing camera, the 160~mm AprilTag marker at the center of the deck of 300~mm
radius, and the wheeled cart that carries the deck at 0.3~m/s.}
\label{fig:setup}
\end{figure}

\section{Results}
\label{sec:results}

The vision-servoed controller landed with a touchdown offset of $46 \pm 18$~mm, against
$187 \pm 92$~mm for the open-loop controller (Table~\ref{tab:results},
Fig.~\ref{fig:offsets}). The reduction is 75\%, or a factor of 4.1. The paired mean
difference is 141~mm with a 95\% confidence interval of [98, 184]~mm
($p = 2 \times 10^{-4}$). The standard deviation of the offset falls by a factor of 5.1,
from 92 to 18~mm.

All 12 vision-servoed landings ended inside the 300~mm deck radius. Nine of the 12
open-loop landings did so, and the remaining three touched down beyond the deck edge.

Descent duration was $14.2 \pm 1.9$~s for the vision-servoed controller and
$9.8 \pm 0.7$~s for the open-loop controller, an increase of 4.4~s, or 45\%.

One vision-servoed trial contained two intervals longer than $\tau_{h}$ in which no marker
detection arrived. The half-amplitude hold was in force during both intervals, detections
resumed in each case, and the trial ended within the deck radius.

\begin{table}[!t]
\caption{Landing performance over $N = 12$ paired trials per controller. Values are mean
$\pm$ standard deviation, and the change column is the paired mean difference with its
95\% confidence interval.}
\label{tab:results}
\centering
\footnotesize
\setlength{\tabcolsep}{3pt}
\begin{tabular}{lrrr}
\toprule
Metric & Open-loop & Vision-servoed & Change \\
\midrule
Touchdown offset (mm) & $187 \pm 92$  & $46 \pm 18$    & $-141$ [$-184$, $-98$] \\
Landings within deck  & 9 / 12        & 12 / 12        & --- \\
Descent duration (s)  & $9.8 \pm 0.7$ & $14.2 \pm 1.9$ & $+4.4$ \\
\bottomrule
\end{tabular}
\end{table}

\begin{figure}[!t]
\centering
% \includegraphics[width=\columnwidth]{fig_offsets.pdf}
\fbox{\parbox[c][22mm][c]{0.92\columnwidth}{\centering\footnotesize
Placeholder: per-trial touchdown offsets for both controllers, paired trials joined by
lines, with the 300~mm deck radius drawn as a reference line.}}
\caption{Touchdown offset for each of the 12 paired landings. Filled circles mark the
open-loop controller and open squares the vision-servoed controller, with a line joining
the two members of each pair. The horizontal reference line is the 300~mm deck radius,
beyond which the vehicle misses the deck.}
\label{fig:offsets}
\end{figure}
